% Notebook-based seed for Lesson 36: no lesson_36.ipynb is present.
% Select the existing bare criticality review and the two-group review that
% matches pages/lesson_36.md; the unrelated control material is omitted.
% Focus source: ../pages/lesson_36.md (Learning Objectives).
\documentclass[compact]{ne630handout}
\handoutsolutionsfalse
% Assigned reading: ../administrivia/schedule-source.psv.
\begin{document}
\Lesson{36}{Two-Group Diffusion and Criticality}{Duderstadt and Hamilton 7.III.B--C; FNRP 7.4.}{Students will be able to write down the equations necessary to model multi-region reactors using the two-group approximation and to solve such equations numerically.}
\begin{handoutflow}

\Sec{1}{Review: Bare-Reactor Criticality}
% Notebook source: ../notebooks/lesson_35.ipynb
% Notebook cell 5 (markdown)
\[
\begin{split}
  B_m^2 \equiv \frac{k_{\infty}-1}{L^2} &= \left ( \frac{\pi}{a} \right )^2 \equiv B_g^2 = \text{geometric buckling}\\
 \rightarrow k_{\infty} &= L^2 \left ( \frac{\pi}{a} \right )^2 + 1 \\
 \rightarrow 1          &= \frac{k_{\infty}}{ 1 + L^2 \left ( \frac{\pi}{a} \right )^2  } = k_{\infty} P_{NL} \\ & = \frac{\text{gains}}{\text{losses from absorption}} \times \frac{\text{losses from absorption}}{\text{losses}} \\
           &= \frac{ \nu\Sigma_f }{ \Sigma_a + D \left ( \frac{\pi}{a} \right )^2  } = \frac{\text{gains}}{\text{losses}} \, .
\end{split}
\]
% End notebook cell 5

% Notebook source: ../notebooks/lesson_35.ipynb
% Notebook cell 6 (markdown)
Each of these statements corresponds to \emph{criticality}, i.e., when losses are exactly balances by gains:

\[
   {\boxed{k = k_{\infty} P_{NL} =   1}} \, .
\]

Just like a real critical system, we get a steady state flux only if there is no source!
% End notebook cell 6

% Notebook source: ../notebooks/lesson_35.ipynb
% Notebook cell 7 (markdown)
But what happens when \(B_m^2 \neq B_g^2\)? As described in Section 7.2, we can still characterize the steady state flux (i.e., solve the NDE) by introducing \(k\) as a new \emph{parameter}. For slab geometry, that looks like

\[
  -D \frac{d^2\phi}{dx^2} + \Sigma_a \phi(x) = \frac{\nu\Sigma_f}{k} \phi(x) \, .
\]

or

\[
 \frac{d^2\phi}{dx^2} + \overbrace{\left ( \frac{k_{\infty}/k-1}{L^2} \right )}^{B^2} \phi(x)  = 0\, .
 \tag{Lewis (7.4), but for slabs}
\]

In other words, \(k\) reduces or increases \(k_{\infty}\) (or, as suggested in the text, \(\nu\)) so that \(B_m^2 = B_g^2\). Our goal, then, is either (1) to determine \(k\) given the materials and geometry or (2) to determine the materials or geometry that leads to a given \(k\) (usually 1!).
% End notebook cell 7

% Notebook source: ../notebooks/L39_f24.ipynb
% Notebook cell 1 (markdown)
\Sec{2}{Review: Two-Group Diffusion}

The two-group equations in slab geometry are \[
 -D_1 \frac{d^2 \phi_1}{dx^2} + \Sigma_{r,1}\phi_1(x)
      = \frac{1}{k_{\text{eff}}} \Big ( \bar{\nu}\Sigma_{f,1}\phi_1(x) + \bar{\nu}\Sigma_{f,2}\phi_2(x) \Big ) \, ,
\]

and

\[
 -D_2 \frac{d^2 \phi_2}{dx^2} + \Sigma_{a,2}\phi_2(x)
      = \Sigma_{s,2\gets 1} \phi_1(x) \, .
\]
% End notebook cell 1

% Notebook source: ../notebooks/L39_f24.ipynb
% Notebook cell 2 (markdown)
What does \(\Sigma_{r,1}\) represent? What are we assuming about \(\Sigma_{s,1\gets 2}\)?
% End notebook cell 2

% Notebook source: ../notebooks/L39_f24.ipynb
% Notebook cell 3 (markdown)
\columnbreak
\Sec{3}{The 1.5-group (aka modified 1-group) Approximation}
% End notebook cell 3

% Notebook source: ../notebooks/L39_f24.ipynb
% Notebook cell 4 (markdown)
Assumptions:

\begin{enumerate}
\def\labelenumi{\arabic{enumi}.}
\item
  Thermal leakage is negligible in the two-group equations, leading to \(\phi_2(x) = (\Sigma_{s_{2\gets 1}} /\Sigma_{a_2})\phi_1(x)\).
\item
  Replaces \(L_1^2\) with \(M^2 = L_1^2 + L_2^2\).
\end{enumerate}
% End notebook cell 4

% Notebook source: ../notebooks/L39_f24.ipynb
% Notebook cell 5 (markdown)
For a two-region reactor (I=fuel, II=reflector), we then have

\[
  \frac{d^2\phi^{\text{I}}_1}{dx^2} +  \frac{  \frac{k_{\infty}}{k}  - 1}{(M^{\text{I}})^2} \phi^{\text{I}}_1(x) = 0 \, ,
\]

and

\[
  \frac{d^2\phi^{\text{II}}_1}{dx^2} - \frac{1}{(M^{\text{II}})^2} \phi^{\text{II}}_1(x) = 0 \, .
\]
% End notebook cell 5

% Notebook source: ../notebooks/L39_f24.ipynb
% Notebook cell 6 (markdown)
\Sec{4}{Full Two-Group Solutions}

Without the 1.5-group approximation, things are a bit more complicated!
% End notebook cell 6

% Notebook source: ../notebooks/L39_f24.ipynb
% Notebook cell 7 (markdown)
The flux in the fuel region is \[
\begin{aligned}
  \phi^{\text{I}}_1(x) =   C_1 \cos(\tilde{B}^{\text{I}}_1 x) + C_2 \cosh(|\tilde{B}^{\text{I}}_2| x)
   \quad \text{and} \quad \\
  \phi^{\text{I}}_1(x) = C_3 \cos(\tilde{B}^{\text{I}}_1 x) + C_4 \cosh(|\tilde{B}^{\text{I}}_2| x)  \, ,
\end{aligned}
\] and \[
\begin{aligned}
  \phi^{\text{II}}_1(x) = C_5 \cosh(|\tilde{B}^{\text{II}}_1| x) + C_6 \sinh(|\tilde{B}^{\text{II}}_2| x)
   \quad \text{and} \quad \\
   \phi^{\text{II}}_2(x) = C_7 \cosh(|\tilde{B}^{\text{II}}_1| x) + C_8 \sinh(|\tilde{B}^{\text{II}}_2| x) \, ,
\end{aligned}
\] where \(\tilde{B}^{i}_j\) are the roots to the quadratic equation for \(B^2\) defined by \[
  \left | \begin{matrix}
    B^2 D^{i}_1 + \Sigma^{i}_{r_1} - \nu\Sigma^{i}_{f_1}/k & -\nu\Sigma_{f_2}^{i}/k \\
    -\Sigma^{i}_{s_{2\gets 1}} & B^2 D^{i}_2 + \Sigma^{i}_{a_2}
   \end{matrix} \right | = 0 \, .
\]

Remember, we can set any of the coefficients \(C_i\) to 1 without changing the solution!
% End notebook cell 7

\end{handoutflow}
\end{document}
